% Experiment Designs poster — NEMI 2026, v1
% Six experiment designs from other fields that MI hasn't adopted.
% Combined framing: designs menu + complementation as centerpiece.
%
%   cd poster && lualatex designs_poster_v1.tex

\documentclass[25pt, a0paper, portrait]{tikzposter}

\usepackage[T1]{fontenc}
\usepackage{amsmath,amssymb}
\usepackage{array}
\usepackage{booktabs}
\usepackage{enumitem}
\usepackage{microtype}

\tikzposterlatexaffectionproofoff

% ── Color palette ───────────────────────────────────────
\definecolor{darktext}{HTML}{2C3E50}
\definecolor{blockbg}{HTML}{FFFFFF}
\definecolor{hookred}{HTML}{C0392B}
\definecolor{designblue}{HTML}{2980B9}
\definecolor{centerpiecegreen}{HTML}{27AE60}

% ── Theme ───────────────────────────────────────────────
\usetheme{Default}

\definecolorstyle{DesignColors}{
  \definecolor{colorOne}{HTML}{2C3E50}
  \definecolor{colorTwo}{HTML}{3498DB}
  \definecolor{colorThree}{HTML}{ECF0F1}
}{
  \colorlet{backgroundcolor}{colorThree}
  \colorlet{framecolor}{colorOne}
  \colorlet{titlefgcolor}{white}
  \colorlet{titlebgcolor}{colorOne}
  \colorlet{blocktitlebgcolor}{designblue}
  \colorlet{blocktitlefgcolor}{white}
  \colorlet{blockbodybgcolor}{blockbg}
  \colorlet{blockbodyfgcolor}{darktext}
  \colorlet{notefgcolor}{darktext}
  \colorlet{notebgcolor}{designblue!12}
  \colorlet{notefrcolor}{designblue}
}
\usecolorstyle{DesignColors}

\defineblockstyle{DesignBlock}{
  titlewidthscale=1.0, bodywidthscale=1.0, titleleft,
  titleoffsetx=0pt, titleoffsety=0pt,
  bodyoffsetx=0pt, bodyoffsety=0pt, bodyverticalshift=0pt,
  roundedcorners=12, linewidth=2pt, titleinnersep=10pt, bodyinnersep=14pt
}{
  \draw[rounded corners=\blockroundedcorners, color=blocktitlebgcolor,
        fill=blockbodybgcolor, line width=\blocklinewidth]
    (blockbody.south west) rectangle (blocktitle.north east);
  \ifBlockHasTitle
    \draw[rounded corners=\blockroundedcorners,
          fill=blocktitlebgcolor, color=blocktitlebgcolor,
          line width=\blocklinewidth]
      (blocktitle.south west) rectangle (blocktitle.north east);
  \fi
}
\useblockstyle{DesignBlock}

\definetitlestyle{DesignTitle}{
  width=\paperwidth, roundedcorners=0, linewidth=0pt, innersep=20pt,
  titletotopverticalspace=15mm, titletoblockverticalspace=25mm
}{
  \draw[fill=titlebgcolor, color=titlebgcolor]
    (\titleposleft, \titleposbottom)
    rectangle (\titleposright, \titlepostop);
}
\usetitlestyle{DesignTitle}

\title{%
  \parbox{0.92\linewidth}{\centering\bfseries
    Experiment Designs We Haven't Run Yet}}
\author{Elliot Tower\quad\texttt{elliot@elliottower.ai}}
\institute{}

\begin{document}

\maketitle[width=\paperwidth]

% ════════════════════════════════════════════════════════
% ROW 1 — THE PITCH + COMPLEMENTATION (centerpiece)
% ════════════════════════════════════════════════════════
\begin{columns}
\column{0.45}

{
\colorlet{blocktitlebgcolor}{hookred}
\block{What's the Sham Control for an Ablation?}{
  Cardiac stenting for stable angina was standard practice for three
  decades---500,000 procedures a year in the US. The first
  sham-controlled trial (ORBITA, 2017) found no significant improvement
  in exercise tolerance versus sham surgery. The intervention was real.
  The causal attribution to the specific intervention had never been
  tested against a matched control.

  \vspace{5mm}
  Ablating a circuit and measuring degradation shows the circuit
  matters. But what's the sham? A size-matched random subnetwork,
  ablated the same way. Without that control, faithfulness scores carry
  the same ambiguity that exercise-tolerance scores carried before
  ORBITA.

  \vspace{5mm}
  Six experiment designs from fields that learned these lessons decades
  ago. Each tests something current methods leave open.

  \vspace{5mm}
  \begin{center}
  \renewcommand{\arraystretch}{1.3}
  {\small
  \begin{tabular}{@{}lll@{}}
    \toprule
    \textbf{Design} & \textbf{Tests} & \textbf{Cost} \\
    \midrule
    Complementation     & Same or different?   & 1 condition \\
    Double dissociation & Specificity          & 1 condition \\
    Dose--response      & Graded vs.\ threshold & Free \\
    Sham control        & Specific vs.\ generic & 1 run \\
    Rescue              & Necessity vs.\ damage & 1 pass \\
    Offset coupling     & Bidirectional link   & Expensive \\
    \bottomrule
  \end{tabular}}
  \end{center}
}
}

\column{0.55}

{
\colorlet{blocktitlebgcolor}{centerpiecegreen}
\block{Complementation: Are These Two Mechanisms or One?}{
  Two head sets both suppress the repeated name in IOI: S-inhibition
  heads and negative name movers. Each degrades performance when
  ablated alone. Are they two independent failure modes, or two names
  for the same functional unit?

  \vspace{5mm}
  \textbf{The test.} Ablate both simultaneously. If the combined
  lesion produces the same deficit as either single lesion, they
  operate through the same pathway---they \emph{complement}. If the
  combined deficit is additive, they are independent mechanisms.

  \vspace{5mm}
  \textbf{Why it matters.} Circuit papers routinely identify multiple
  component classes (``name movers,'' ``S-inhibition heads,'' ``backup
  name movers'') without testing whether these classes are functionally
  distinct. The complementation test answers that question with one
  extra ablation condition.

  \vspace{5mm}
  \textbf{What it looks like.}

  \vspace{3mm}
  \begin{center}
  \renewcommand{\arraystretch}{1.3}
  {\large
  \begin{tabular}{@{}lcl@{}}
    \toprule
    \textbf{Ablation} & \textbf{Deficit} & \textbf{Interpretation} \\
    \midrule
    S-inhibition only     & $\Delta_1$ & \\
    Neg.\ name movers only & $\Delta_2$ & \\
    \midrule
    Both (if $\approx \max$)  & $\approx \Delta_1$ & Same pathway \\
    Both (if $\approx$ sum)   & $\Delta_1 + \Delta_2$ & Independent \\
    \bottomrule
  \end{tabular}}
  \end{center}

  \vspace{5mm}
  Across sixteen published mechanistic claims, this test has been
  applied \textbf{zero times}.

  \vspace{4mm}
  \textbf{Cost:} one extra ablation condition.

  \vspace{3mm}
  {\footnotesize Genetics. Benzer's cis-trans test (1955), standard in
  functional genomics.}
}
}

\end{columns}

% ════════════════════════════════════════════════════════
% ROW 2 — THREE DESIGNS
% ════════════════════════════════════════════════════════
\begin{columns}
\column{0.33}

\block{Double Dissociation}{
  Two interventions that cross: A breaks function X and spares
  function Y; B breaks Y and spares X. One arm alone shows
  involvement. The crossing shows specificity.

  \vspace{4mm}
  \textbf{In MI:} patching a belief-tracking subspace in Llama-3-70B
  yields IIA~$= 1.0$ on belief-framed questions and $0.0$ on
  reality-framed questions---about identical underlying facts. The
  subspace responds to the epistemic framing, not the information.

  \vspace{4mm}
  Across sixteen published claims, exactly \textbf{one} meets this
  standard---by different authors, three years after the original
  paper.

  \vspace{4mm}
  \textbf{Cost:} one extra condition.

  \vspace{3mm}
  {\footnotesize Neuropsychology. Standard since Teuber (1955).}
}

\column{0.34}

\block{Dose--Response}{
  Partial intervention should produce partial effect. A step function
  means you measured a threshold, not a mechanism.

  \vspace{4mm}
  \textbf{In MI:} sweep ablation strength from 0\% to 100\% instead
  of reporting one on/off number. The shape of the curve distinguishes
  specific mechanisms (graded, monotonic) from nonspecific disruption
  (all-or-nothing collapse at some threshold).

  \vspace{4mm}
  Most published evaluations report a single ablation point. The curve
  between zero and complete ablation carries more information than
  either endpoint.

  \vspace{4mm}
  \textbf{Cost:} free---you already have the interpolation.

  \vspace{3mm}
  {\footnotesize Pharmacology. Dose-response curves, receptor theory.}
}

\column{0.33}

\block{Sham Control}{
  Not ``no intervention'' but ``the same procedure with the active
  ingredient removed.''

  \vspace{4mm}
  \textbf{In MI:} ablate a size-matched random subnetwork---same count,
  same layer distribution---and compare. If the identified circuit and
  the random circuit degrade performance equally, the specificity claim
  is unfounded.

  \vspace{4mm}
  ORBITA showed that three decades of clinical confidence rested on
  comparison against no intervention rather than against a matched
  control. The same gap exists in most circuit evaluations.

  \vspace{4mm}
  \textbf{Cost:} one extra run.

  \vspace{3mm}
  {\footnotesize Clinical trials. Sham surgery, placebo controls.}
}

\end{columns}

% ════════════════════════════════════════════════════════
% ROW 3 — TWO MORE DESIGNS + TAKEAWAY
% ════════════════════════════════════════════════════════
\begin{columns}
\column{0.33}

\block{Rescue / Reversibility}{
  Remove a component, observe the deficit, then restore it. If
  function does not recover, the ablation caused collateral damage
  beyond what the mechanism hypothesis predicts.

  \vspace{4mm}
  \textbf{In MI:} re-inject the clean activation after ablation and
  measure recovery. Rescue separates ``this component is necessary''
  from ``removing this component broke something downstream.''

  \vspace{4mm}
  Activation patching already contains the rescue logic. Almost no
  published evaluation reports the recovery rate alongside the
  degradation rate.

  \vspace{4mm}
  \textbf{Cost:} one forward pass.

  \vspace{3mm}
  {\footnotesize Genetics. Rescue and complementation experiments.}
}

\column{0.34}

\block{Offset Coupling}{
  A mechanism should disappear when its capability disappears.
  Onset coupling---the circuit forms when the capability appears
  during training---is reported for induction heads. The partner
  test is not.

  \vspace{4mm}
  \textbf{In MI:} unlearn or fine-tune away a capability, then check
  whether the circuit degrades in tandem. If the circuit persists after
  the capability is gone, the causal link runs in one direction only.

  \vspace{4mm}
  Onset coupling is suggestive. Offset coupling is confirmatory.
  Together they establish a bidirectional association between mechanism
  and capability.

  \vspace{4mm}
  \textbf{Cost:} expensive (fine-tuning or unlearning required).

  \vspace{3mm}
  {\footnotesize Developmental neuroscience. Lesion-recovery studies.}
}

\column{0.33}

{
\colorlet{blocktitlebgcolor}{colorOne}
\block{Takeaway}{
  {\Large\bfseries Each of these designs tests something that current
  methods leave open. Most cost one extra condition.}

  \vspace{5mm}
  \begin{itemize}[leftmargin=*, itemsep=3mm]
    \item \textbf{3 of 6} designs cost one extra condition
    \item \textbf{1 of 6} is free
    \item \textbf{1 of 16} published claims meets double dissociation
    \item \textbf{0 of 16} meets complementation
    \item Combined methodological heritage: 250+ years
  \end{itemize}

  \vspace{5mm}
  {\small A claim is as strong as its weakest validity dimension. These
  designs close the dimensions that current methods leave untested.}

  \vspace{5mm}
  \begin{center}
  \texttt{elliot@elliottower.ai}

  \vspace{2mm}
  {\small Framework:} \textit{Mechanistic Validity}

  {\small 36 criteria, 16 claims audited}
  \end{center}
}
}

\end{columns}

\end{document}
